\documentclass[12pt,a4paper]{article}
\usepackage[utf8]{inputenc}
\usepackage[T1]{fontenc}
% \usepackage[brazilian]{babel} % Removido para evitar erro de pacote não instalado
\renewcommand{\contentsname}{Sumário}
\usepackage{geometry}
\geometry{a4paper, left=3cm, top=3cm, right=2cm, bottom=2cm}
\usepackage{setspace}
\usepackage{indentfirst}
\usepackage{listings}
\usepackage{xcolor}
\usepackage{hyperref}

\setstretch{1.5} % ABNT: espaçamento 1.5

% Configuração para listagens de código
\definecolor{codegreen}{rgb}{0,0.6,0}
\definecolor{codegray}{rgb}{0.5,0.5,0.5}
\definecolor{codepurple}{rgb}{0.58,0,0.82}
\definecolor{backcolour}{rgb}{0.95,0.95,0.92}

\lstdefinestyle{mystyle}{
    backgroundcolor=\color{backcolour},   
    commentstyle=\color{codegreen},
    keywordstyle=\color{blue},
    numberstyle=\tiny\color{codegray},
    stringstyle=\color{codepurple},
    basicstyle=\ttfamily\footnotesize,
    breakatwhitespace=false,         
    breaklines=true,                 
    captionpos=b,                    
    keepspaces=true,                 
    numbers=left,                    
    numbersep=5pt,                  
    showspaces=false,                
    showstringspaces=false,
    showtabs=false,                  
    tabsize=2,
    extendedchars=true,
    literate={á}{{\'a}}1 {ã}{{\~a}}1 {é}{{\'e}}1 {í}{{\'i}}1 {ó}{{\'o}}1 {õ}{{\~o}}1 {ú}{{\'u}}1 {ç}{{\c{c}}}1 {Á}{{\'A}}1 {Ã}{{\~A}}1 {É}{{\'E}}1 {Í}{{\'I}}1 {Ó}{{\'O}}1 {Õ}{{\~O}}1 {Ú}{{\'U}}1 {Ç}{{\c{C}}}1
}
\lstset{style=mystyle}

\begin{document}

% =================== CAPA ===================
\begin{titlepage}
    \begin{center}
        {\Large \textsc{Universidade de Brasília -- UnB}}\\[0.2cm]
        {\large Departamento de Ciência da Computação -- CIC}\\[0.2cm]
        {\large Disciplina: Redes de Computadores}\\[3cm]

        {\large \textbf{Gabriel Gonçalves Caldo} -- 231034627}\\
        {\large \textbf{Mario Augusto Vieira dos Santos} -- 231035778}\\
        {\large \textbf{Daniel Rodrigues de Abreu} -- 241038540}\\[0.5cm]
        {\large \textbf{Grupo 2}}\\[4cm]

        \hrule
        \vspace{0.8cm}
        {\huge \textbf{RELATÓRIO TÉCNICO}}\\[0.4cm]
        {\LARGE \textit{Projeto de Chat P2P (Versão Revisada)}}\\[0.8cm]
        \hrule
        \vfill

        {\large \textbf{Brasília, DF}}\\
        {\large \textbf{\the\year}}
    \end{center}
\end{titlepage}

% =================== FOLHA DE ROSTO ===================
\begin{titlepage}
    \begin{center}
        {\large \textbf{Gabriel Gonçalves Caldo} -- 231034627}\\
        {\large \textbf{Mario Augusto Vieira dos Santos} -- 231035778}\\
        {\large \textbf{Daniel Rodrigues de Abreu} -- 241038540}\\[0.5cm]
        {\large \textbf{Grupo 2}}\\[4cm]

        \hrule
        \vspace{0.8cm}
        {\huge \textbf{RELATÓRIO TÉCNICO}}\\[0.4cm]
        {\LARGE \textit{Projeto de Chat P2P (Versão Revisada)}}\\[0.8cm]
        \hrule
        \vspace{3cm}

        \begin{flushright}
        \begin{minipage}{0.5\textwidth}
            \normalsize
            Relatório técnico descritivo da arquitetura e protocolos de rede, apresentado como requisito parcial de avaliação para a disciplina de Redes de Computadores da Universidade de Brasília (UnB).
        \end{minipage}
        \end{flushright}
        \vfill

        {\large \textbf{Brasília, DF}}\\
        {\large \textbf{\the\year}}
    \end{center}
\end{titlepage}

\newpage
% =================== SUMÁRIO ===================
\tableofcontents
\newpage

% =================== CONTEÚDO ===================

\section{Introdução}
Este relatório técnico descreve a implementação do cliente de Chat P2P desenvolvido pelo Grupo 2 para a disciplina de Redes de Computadores da Universidade de Brasília (UnB). O sistema é capaz de se registrar e buscar pares em um servidor Rendezvous central, além de estabelecer conexões diretas TCP (\textit{Peer-to-Peer}) persistentes com outros clientes para a troca de mensagens, seguindo os requisitos estipulados na especificação do projeto. A comunicação ocorre 100\% em formato P2P após a descoberta de pares, isto é, sem saltos intermediários (\textit{relay}) e com a ausência de um ponto centralizado de troca de mensagens.

\section{Arquitetura da Solução}
A arquitetura do software foi concebida aproveitando ao máximo a biblioteca nativa assíncrona do Python (\texttt{asyncio}). Isto permite o tratamento concorrente de múltiplas conexões de rede numa \textit{thread} única, conferindo eficiência e simplicidade no controle de estado.


Abaixo, apresentamos o diagrama estrutural de como os módulos principais se organizam e interagem de forma descentralizada:

\begin{lstlisting}[basicstyle=\ttfamily\small, numbers=none]
+-----------------------------------------------------+
|                      main.py                        |
|  (Ponto de entrada, inicializa loop e dependências) |
+--------------------------+--------------------------+
                           |
          +----------------v----------------+
          |             cli.py              |
          |  (Interface interativa do CLI)  |
          +----------------+----------------+
                           |
+--------------------------v--------------------------+
|                 message_router.py                   |
|       (Roteia as mensagens: SEND, PUB, ACK)         |
+-----+--------------------+--------------------+-----+
      |                    |                    |
+-----v------+       +-----v------+       +-----v------+
|peer_server |       |peer_conn.  |       |rendezvous_ |
| (inbound)  |       | (outbound) |       | connection |
+------------+       +------------+       +------------+
      |                    |                    |
+-----v--------------------v--------------------v-----+
|                   peer_table.py                     |
|        (Armazena o estado e tabela da rede)         |
+-----------------------------------------------------+
\end{lstlisting}

Nossa implementação divide a lógica em diversos módulos desacoplados:
\begin{itemize}
    \item \textbf{\texttt{main.py}}: Arquivo principal. Carrega as definições, registra no servidor central via \texttt{rendezvous\_connection.py} e inicia o ciclo do \textit{event loop} do \texttt{asyncio}, englobando interface com usuário e servidor.
    \item \textbf{\texttt{peer\_server.py}}: Define o servidor local para escuta de conexões \textit{inbound} (recebidas) de outros pares P2P.
    \item \textbf{\texttt{peer\_connection.py}}: Implementa o manipulador de cada conexão TCP bidirecional (\textit{inbound} ou \textit{outbound}). Ele abstrai a leitura delimitada por \texttt{\textbackslash n}, lida com pacotes \textit{keep-alive} (\texttt{PING/PONG}), negocia a sessão inicial com \texttt{HELLO/HELLO\_OK} e despacha as mensagens de chat recebidas.
    \item \textbf{\texttt{rendezvous\_connection.py}}: Concentra a requisição sobre \textit{sockets} para interagir com o servidor de Rendezvous, tratando comandos de \texttt{REGISTER}, \texttt{DISCOVER} e \texttt{UNREGISTER}.
    \item \textbf{\texttt{message\_router.py}}: O \textit{Router} é responsável pela difusão e roteamento em tempo real das saídas dos chats. Se for uma mensagem \textit{unicast} (direcionada a apenas um par, tipo \texttt{SEND}), ele busca a conexão específica. Se for \textit{broadcast} (tipo \texttt{PUB}), despacha a todos os pares válidos.
    \item \textbf{\texttt{peer\_table.py}}: Funciona como a tabela de vizinhos ou estado (\textit{State}). Ele armazena metadados de outros pares registrados na mesma sala (\textit{namespace}) e controla a saúde dos enlaces, marcando pares obsoletos como inativos.
    \item \textbf{\texttt{cli.py}}: Fornece a interface visual baseada em linha de comando interativa para acesso humano. Transmite chamadas dos comandos (ex: \texttt{/msg}, \texttt{/pub}, \texttt{/conn}) para as outras camadas.
\end{itemize}

O cliente lê dados seriais formatados estritamente em JSON de UTF-8. Assim, a comunicação possui tamanho restrito e delimitador forte para a quebra de 32KB em buffer TCP, conforme exigido.

\section{Protocolos Implementados com Exemplos de Funcionamento}

O projeto concretiza diferentes interações ao nível de aplicação na rede, conforme o seguinte mapeamento. Todos os comandos JSON terminam em um caractere de nova linha (\texttt{\textbackslash n}).

\subsection{Estabelecimento e Sessão Inicial (\texttt{HELLO} / \texttt{HELLO\_OK})}
Todo enlace começa com o reconhecimento de identidade do \textit{peer}.
\begin{itemize}
    \item \textbf{Cenário:} Alice conecta-se fisicamente no IP do Bob e emite um HELLO. Bob responde informando que aceita o fluxo de dados através do HELLO\_OK.
    \item \textbf{Exemplo Requisição (Alice $\rightarrow$ Bob):}
\begin{lstlisting}[numbers=none]
{
  "type": "HELLO", 
  "peer_id": "alice@UnB", 
  "version": "1.0", 
  "features": [],
  "ttl": 1
}
\end{lstlisting}
    \item \textbf{Exemplo Resposta (Bob $\rightarrow$ Alice):}
\begin{lstlisting}[numbers=none]
{
  "type": "HELLO_OK", 
  "peer_id": "bob@UnB"
}
\end{lstlisting}
\end{itemize}

\subsection{Keep-Alive (\texttt{PING} / \texttt{PONG})}
Para mensurar RTT (Tempo de Ida e Volta) e evidenciar que a conexão TCP não foi morta por \textit{middleboxes} como NATs e \textit{firewalls}, enviamos ping periódico:
\begin{itemize}
    \item \textbf{Exemplo PING:}
\begin{lstlisting}[numbers=none]
{
  "type": "PING", 
  "peer_id": "alice@UnB", 
  "version": "1.0", 
  "features": [],
  "ttl": 1
}
\end{lstlisting}
    \item \textbf{Exemplo PONG:}
\begin{lstlisting}[numbers=none]
{
  "type": "PONG"
}
\end{lstlisting}
\end{itemize}

\subsection{Mensagens de Chat e Encerramento}
As mensagens usam a marcação \texttt{SEND} (Unicast), \texttt{ACK} (Recibo de Recebimento) e \texttt{PUB} (Broadcast). Ao fim de tudo, \texttt{BYE} e \texttt{BYE\_OK} gerenciam o fechamento.
\begin{itemize}
    \item \textbf{Exemplo de SEND (Unicast com exigência de recibo):}
\begin{lstlisting}[numbers=none]
{
  "type": "SEND", 
  "msg_id": "550e8400-e29b-41d4-a716-446655440000",
  "src": "alice@UnB", 
  "dst": "bob@UnB", 
  "payload": "Olá, tudo bem?", 
  "require_ack": true, 
  "ttl": 1
}
\end{lstlisting}
    \item \textbf{Exemplo de ACK correspondente:}
\begin{lstlisting}[numbers=none]
{
  "type": "ACK", 
  "msg_id": "uuid-da-mensagem-anterior", 
  "timestamp": "2025-10-27T10:00:01Z", 
  "ttl": 1
}
\end{lstlisting}
    \item \textbf{Exemplo de BYE (Encerramento do enlace TCP):}
\begin{lstlisting}[numbers=none]
{
  "type": "BYE", 
  "reason": "CLI Fechando a conexao"
}
\end{lstlisting}
\end{itemize}

\section{Instruções Detalhadas para Execução}

Esta subseção responde às perguntas contidas na documentação técnica para viabilizar a implantação local pelos usuários ou avaliadores.

\subsection{Requisitos Mínimos para Executar o Programa em Linux}
Sim, \textbf{é obrigatório que o programa rode no Linux}. Os requisitos mínimos são:
\begin{enumerate}
    \item Sistema Operacional derivado do Kernel Linux (Ubuntu 20.04+, Debian, Fedora, Arch, etc.).
    \item Interpretador \textbf{Python versão 3.8 ou superior} nativamente instalado.
    \item Idealmente, firewalls e proteções semelhantes \textbf{desligados} para evitar problemas na conexão com o servidor Rendezvous e outros peers.
    \item Privilégios regulares de usuário para realizar a abertura (\textit{binding}) de \textit{sockets} nas portas do intervalo dinâmico não privilegiado (portas de 1024 a 65535, configuradas no projeto).
\end{enumerate}

\subsection{Compilação do Código}
Como o ecossistema é alicerçado sob a linguagem \textbf{Python}, não é feito uso de nenhuma linguagem \textit{Ahead-of-Time} (AoT) compilada, como C ou Java. Assim, \textbf{não deve ser feita nenhuma compilação prévia do código}. O código é interpretado dinamicamente durante o tempo de execução através do \textit{bytecode} do interpretador \texttt{python3}.

\subsection{Bibliotecas Necessárias e Como Instaladas}
O projeto arquitetural atendeu rigorosamente à solicitação de evitar dependências complexas e \textit{frameworks} exógenos. Desta forma, \textbf{apenas a Biblioteca Padrão do Python (Standard Library)} foi utilizada. Bibliotecas como \texttt{asyncio}, \texttt{socket}, \texttt{json}, \texttt{logging} e \texttt{datetime} estão nativamente embarcadas no Python.
\textbf{Não existem bibliotecas extras exigidas para o uso do programa}. O documento \texttt{requirements.txt} inclui estritamente artefatos suplementares de formatação recomendados (como o \texttt{black} e o \texttt{flake8}) unicamente voltados a ambiente de desenvolvimento (\textit{dev-environment}), cuja instalação é totalmente omitível.
Caso mesmo assim haja interesse nestas ferramentas de estilo:
\begin{lstlisting}[language=bash, numbers=none]
$ python3 -m venv .venv
$ source .venv/bin/activate
$ pip install -r requirements.txt
\end{lstlisting}

\subsection{Passos Necessários para Executar o Programa}
A cadeia de passos sequenciais que guiam o sistema ao pleno funcionamento no terminal Linux é:
\begin{enumerate}
    \item Faça a abertura de um terminal e navegue via \texttt{cd} até o diretório da raiz do código fonte que contém os arquivos Python:
\begin{lstlisting}[language=bash, numbers=none]
$ cd /caminho/ate/Redes_Computadores/
\end{lstlisting}
    \item Opcionalmente, pode-se editar o arquivo de metadados padrão \texttt{config.json} para adequar o \texttt{namespace}, porta \textit{inbound} e identidade do \textit{peer}.
    \item Inicialize a aplicação chamando diretamente o script de partida:
\begin{lstlisting}[language=bash, numbers=none]
$ python3 main.py
\end{lstlisting}
    \item O \textit{CLI} interativo abrirá o canal após notificar seu registro com sucesso. Você pode, agora, explorar os comandos:
        \begin{itemize}
            \item \texttt{/peers *} para invocar uma verificação global da mesa Rendezvous.
            \item \texttt{/conn} para analisar ativamente com quem seu hospedeiro está amarrado no protocolo TCP local.
            \item \texttt{/msg [peer\_id] [mensagem]} para realizar o envio \textit{unicast}.
            \item \texttt{/pub * [mensagem]} para lançar pacotes \textit{broadcast}.
        \end{itemize}
    \item Digite \texttt{/quit} quando terminar de usar, para assim desregistrar-se no servidor e finalizar todas as sessões.
\end{enumerate}

\end{document}
